% =====================================================================
%  Glossario  (scansione p-70 .. p-71)
%  STATO: completo (voci da "Cent" a "Varietà di temperamento mesotonico"
%  più la tabella delle proporzioni matematiche/intervalli greci).
% =====================================================================
\chapter*{Glossario}
\addcontentsline{toc}{chapter}{Glossario}
\markboth{Glossario}{Glossario}

\voce{Cent}{Unità di misura dell'intervallo. La centesima parte di un semitono
temperato, con rapporto $\sqrt[1200]{2}$.}

\voce{Circolo delle quinte}{L'arrangiamento di note in un sistema chiuso di
quinte, come Do, Sol, Re, La, Mi, ecc.}

\voce{Comma}{Errore di intonazione, come nell'intervallo Si$\sharp$--Do nel
sistema pitagorico. Vedi \emph{Comma ditonico} e \emph{Comma sintonico}.}

\voce{Comma ditonico}{Intervallo tra 12 quinte e 7 ottave. Il suo rapporto è
$531441/524288$ ed è considerato convenzionalmente di 23,5 cents.}

\voce{Comma pitagorico}{Vedi \emph{Comma ditonico}.}

\voce{Comma sintonico}{Intervallo tra una terza maggiore naturale ($\rapp{5}{4}$)
e la terza pitagorica ($\rapp{81}{64}$). Il suo rapporto è $81/80$ ed è
considerato convenzionalmente di 21,5 cents.}

\voce{Comma tolematico}{Vedi \emph{Comma sintonico}.}

\voce{Diesis}{L'intervallo (circa $1/5$ di tono) tra due note enarmonicamente
equivalenti, come La$\flat$ e Sol$\sharp$, nell'intonazione naturale o nel
temperamento mesotonico. Il suo rapporto è $128/125$ di circa 41 cents.}

\voce{Divisione aritmetica}{Divisione uguale della differenza tra due quantità
tale che il risultato formi una progressione aritmetica; ad es.: $9:8:7:6$.}

\voce{Divisione geometrica}{Divisione proporzionale di due quantità, tale che il
risultato formi una progressione geometrica; ad es.: $27:18:12:8$.}

\voce{Divisione multipla}{Divisione dell'ottava in più di 12 parti, uguali o
ineguali.}

\voce{Epimorio}{Propriamente \enquote{una parte aggiunta}. È il termine usato
dai Greci per definire intervalli superparticolari. Vedi \emph{Rapporto
superparticolare}.}

\voce{Intonazione naturale}{Sistema d'intonazione basato sull'ottava
($\rapp{2}{1}$), la quinta pura ($\rapp{3}{2}$) e la terza maggiore pura
($\rapp{5}{4}$).}

\voce{Intonazione pitagorica}{Sistema d'intonazione basato sull'ottava
($\rapp{2}{1}$) e la quinta pura ($\rapp{3}{2}$).}

\voce{Monocordo}{Corda tesa su una base di legno sulla quale sono indicate le
lunghezze della corda per alcuni sistemi d'intonazione; diagramma che illustra
tali lunghezze; indicazioni per costruire tale diagramma.}

\voce{Quinta del lupo}{La quinta dissonante, usualmente Sol\#--Mi$\flat$ (notata
come sesta diminuita), presente in qualsiasi temperamento ineguale (tale che
sia di 737 cents anziché 702 cents).}

\voce{Rapporto superparticolare}{Un rapporto in cui l'antecedente eccede il
conseguente di 1, ad es. $\rapp{5}{4}$. Vedi \emph{Sesqui}.}

\voce{Schisma}{Differenza tra il comma sintonico e il comma ditonico, con
rapporto $32805/32768$, approssimativamente 2 cents.}

\voce{Sesqui}{Prefisso usato per designare un rapporto superparticolare, ad es.
la sesquiterza ($\rapp{4}{3}$).}

\voce{Sistema chiuso}{Un temperamento regolare in cui la nota iniziale si
ripresenta nuovamente.}

\voce{Sistema di intonazione}{Un sistema in cui tutti gli intervalli possono
essere espressi da numeri razionali.}

\voce{Temperamento equabile}{Divisione dell'ottava in un numero uguale di
parti, specificatamente in 12 semitoni, ciascuno dei quali ha rapporto
$\sqrt[12]{2}$.}

\voce{Temperamento ineguale}{Qualsiasi temperamento diverso da quello equabile,
particolarmente quello mesotonico o altre sue varietà.}

\voce{Temperamento mesotonico}{Propriamente un sistema di intonazione con
quinte abbassate e terze maggiori pure ($\rapp{5}{4}$). Vedi anche
\emph{Varietà di temperamento mesotonico}.}

\voce{Temperamento regolare}{Temperamento in cui tutte le quinte eccetto una
sono delle medesime dimensioni, come ad es. il sistema pitagorico o il
temperamento mesotonico (il temperamento equabile, con tutte le quinte uguali,
è anch'esso un temperamento regolare, così pure i sistemi chiusi di divisione
multipla).}

\voce{Temperare}{Variare leggermente l'altezza; specificatamente di una quinta
temperata.}

\voce{Varietà di temperamento mesotonico}{Temperamenti regolari basati sullo
stesso principio del temperamento mesotonico, con quinte abbassate e terze
alterate.}

\bigskip
\begin{center}
\textbf{Proporzioni matematiche e rispettivi intervalli musicali}\\[4pt]
\begin{tabular}{@{}lll@{}}
\toprule
Rapporto & Nome latino & Intervallo greco \\
\midrule
$\rapp{2}{1}$ & Dupla            & Diapason \\
$\rapp{3}{1}$ & Tripla           & Diapason Diapente \\
$\rapp{4}{1}$ & Quadrupla        & Bis Diapason \\
$\rapp{3}{2}$ & Sesquialtera     & Diapente \\
$\rapp{4}{3}$ & Sesquiterza      & Diatessaron \\
$\rapp{5}{4}$ & Sesquiquarta     & Diatonus Semitonus \\
$\rapp{8}{3}$ & Dupla Superbipartiens & Diapason Diatessaron \\
$\rapp{9}{8}$ & Sesquiottava     & Tonus \\
\bottomrule
\end{tabular}
\end{center}
